\begin{abstract}
    \ji{
    World action models (WAMs) have emerged as a promising approach to robotic control by jointly modeling future visual dynamics and actions. 
    However, the dense spatiotemporal representations required for video generation introduce substantial computational overhead, even though not all visual information is equally relevant to action prediction. 
    In this paper, we propose \method, an efficient WAM framework that allocates computation according to the action relevance of future visual representations. 
    Our key idea is to preserve fine-grained spatiotemporal information where it matters for action generation, while representing less relevant regions at coarser resolutions. 
    To this end, we introduce a lightweight action-aware router trained with gradient-based supervision from the action decoder. 
    The learned relevance guides spatial token merging, preserving the resolution of action-critical regions, while adaptive temporal aggregation jointly considers action relevance and feature dynamics to capture their temporal evolution. 
    We further employ velocity distillation to adapt the action decoder to compressed representations without updating the video generation backbone.
    By selectively allocating spatial and temporal resolution, \method reduces the number of future tokens by up to 90\% in the specified configuration.
    [Quantitative results demonstrating action performance and inference efficiency.]
    }
% World action models (WAMs) build robot policies on pretrained video generation models, so that the action decoder inherits the visual and dynamic knowledge of the video backbone.
% The backbone, however, processes a dense spatio-temporal token grid whose layout is set by video generation, while the action decoder relies on only part of it;
% the rest is redundant for control but is still computed and passed to the decoder.
% Generic token reduction does not resolve this mismatch:
% visual similarity is a poor proxy for action relevance, and deleting tokens degrades the policy sharply when few tokens are kept.
% We propose \method, an action-aware token allocation that lets the action, rather than visual content, decide where and at which resolution the video backbone computes.
% A lightweight router, supervised only by how each token affects the predicted action, scores the future-frame tokens.
\end{abstract}
